\documentclass[11pt]{article}
\usepackage{microtype,amsmath,amssymb,amsthm,mathtools,booktabs,array}
\usepackage[margin=1in]{geometry}
\usepackage[colorlinks=true,linkcolor=blue!45!black,citecolor=blue!45!black,urlcolor=blue!45!black]{hyperref}
\usepackage{xcolor}
\hypersetup{pdftitle={Gaussian Marker Reduction with a Constant Quadratic Fibre},pdfsubject={Gaussian marker reduction and coefficient specialization}}
\newtheorem{theorem}{Theorem}[section]
\newtheorem{proposition}[theorem]{Proposition}
\newtheorem{lemma}[theorem]{Lemma}
\newtheorem{corollary}[theorem]{Corollary}
\theoremstyle{definition}\newtheorem{definition}[theorem]{Definition}
\theoremstyle{remark}\newtheorem{remark}[theorem]{Remark}
\newcommand{\E}{\mathbb E}
\newcommand{\C}{\mathbb C}
\newcommand{\tr}{\operatorname{tr}}
\newcommand{\wt}{\operatorname{wt}}
\newcommand{\G}{\mathsf G}
\DeclareMathOperator{\Span}{span}
\title{Gaussian marker reduction\\with a constant quadratic fibre}
\author{Roy van Rijn}\date{17 November 2026: unpublished manuscript}
\begin{document}\maketitle

\begin{abstract}
Over every characteristic-zero field, a Gaussian phase that is quadratic
in one fibre variable with constant quadratic coefficient has every
polynomially marked moment series generated by base-variable markers,
time derivatives and polynomial time coefficients. Over a general $\mathbb Q$-algebra, we
prove a coefficient-content theorem for the polynomial source and actual
marker generation for nilpotent coefficient content. We also give a
sufficient mixed-moment degree bound. For the family $Y^2+bY$, we compute
the exact actual-image defect: it is a point module with coefficient
module $R/(bR+\Gamma_b(R))$, where $\Gamma_b(R)$ consists of elements
killed by a power of $b$. This gives a necessary and sufficient
generation criterion in that family and shows why generation after
every field specialization and every finite nilpotent quotient need
not descend to the original coefficient ring.
\end{abstract}

\section{A reduction of the original markers}

Let $U=(U_1,\ldots,U_r)$ and $Y$ be independent standard Gaussian
coordinates and let
\[
 P=aY^2+B(U)Y+g(U),\qquad a\text{ constant}.
\]
The degrees of $B,g$ and the number of base variables are arbitrary.
A polynomial marker is a fixed polynomial multiplier $Q$. Use the normalization
\[
 \Phi_Q(t)=\sum_{m\ge0}\E[QP^m]\frac{t^m}{m!}.
\]
The Gaussian functional is defined algebraically by
$\E[U_i^{2j}]=(2j-1)!!$, $\E[U_i^{2j+1}]=0$, independence,
and the same rules for $Y$, with $\E[1]=1$. These rules also define it
over every commutative $\mathbb Q$-algebra. No convergence of an
exponential integral is assumed.

Call $P$ \emph{pure} if $\Phi_1=1$, equivalently
$\E[P^m]=0$ for every $m\ge1$. Call it \emph{safe} if every
$\Phi_Q$ is a polynomial in $t$, equivalently, for each fixed $Q$,
$\E[QP^m]=0$ for all sufficiently large $m$. A base marker depends
only on $U$.
For a characteristic-zero field $K$, let
$D_K=K[t]\langle\partial_t\rangle$ act on formal time series by
multiplication by $t$ and differentiation. The \emph{actual marked
module} is the $D_K$-submodule
$V=\sum_Q D_K\Phi_Q\subset K[[t]]$. Expressions in this module
are finite sums; no completion of the operator ring is used.

\begin{theorem}\label{thm:actual}
Over every characteristic-zero field $K$, the actual marked module $V$
is generated by the periods $\Phi_Q$ with $Q\in K[U]$.
In particular $P$ is safe if and only if all these base-marker periods
are polynomial. Purity is not needed for this equivalence.

Every marker has a finite expression in time derivatives of base-marker
periods, with polynomial time coefficients and explicit spatial
integration-by-parts witnesses.
\end{theorem}

For two base variables this locates any failure in a base marker, but
does not prove that purity implies safety. Conditional integration leaves
a parameter-dependent binary phase, which the fixed-phase binary theorem
does not cover.

The proof combines two recurrences for even fibre powers and uses
spatial differentiation and the time commutator for odd powers.
The resulting coefficient-content theorem gives finite witnesses and
a degree transfer. A finite parity argument also handles nilpotent
coefficient content. We then compute the exact specialization defect
for $P_b=Y^2+bY$, which explains why fieldwise generation need not
descend to a general coefficient ring.

We prove a stronger coefficient-ring statement, which identifies
the specialization issue precisely. Let $R$ be any commutative
$\mathbb Q$-algebra, without reducedness or Noetherian assumptions.
Put
\begin{align*}
 A&=R[t,U,Y],\\
 D_i&=\partial_{U_i}+t(B_iY+g_i)-U_i,
 &D_Y&=\partial_Y-(1-2at)Y+tB,\\
 M&=A\big/\bigl(\sum_iD_iA+D_YA\bigr),
 &\nabla&=\partial_t+P.
\end{align*}
This is a vector-space or $R[t]$-module quotient, not a quotient
by the ordinary Jacobian ideal. The operators $D_i,D_Y$ commute
with $\nabla$, so it is a module over $D_R=R[t]\langle\partial_t\rangle$.
Write $[Q]$ for a class, let $N$ be the $D_R$-submodule generated by
all $[Q]$ with $Q\in R[U]$, and let $I_B\subset R$ be the ideal
of all coefficients of $B$.

\begin{theorem}[Ambient certificates and coefficient content]
\label{thm:content}
Every even fibre class belongs to $N$, and
\[
 I_B(M/N)=0.
\]
These statements have universal polynomial operator certificates.
If $I_B=R$, then $M=N$ over the ring $R$ itself. If $B=0$
identically, the actual Gaussian image is still generated by base
markers, because every odd fibre period is zero by parity.
An ambient odd point-torsion module may remain when $B=0$;
it must not be mistaken for an actual Gaussian image.
\end{theorem}

Over a field, a nonzero $B$ has $I_B=K$, proving
Theorem~\ref{thm:actual}; parity treats $B=0$. Over a general ring,
Corollary~\ref{cor:nilpotentcontent} gives actual generation when
$I_B$ is nilpotent, even if $B\ne0$. Section~\ref{sec:specialization}
shows that an actual-image defect can persist for other nonunit content.

\section{Two polynomial certificates for every even marker}

Set
\[
 s=1-2at,\qquad u=1-at,\qquad 2u-s=1,
 \qquad E_j(Q)=[QY^{2j}].
\]
Here $Q\in R[t,U]$ may depend polynomially on time.
Fibre integration by parts gives, for every $m\ge0$,
\begin{equation}\label{eq:conditional}
 s^m[QY^m]=C_m(Q):=
 \sum_{k=0}^{\lfloor m/2\rfloor}
 \frac{m!}{2^k k!(m-2k)!}\,
       t^{m-2k}s^k[QB^{m-2k}].
\end{equation}
This is a polynomial identity in $M$; no inverse of $s$ is used.
To retain its exact spatial witness, define $W_0=0$, $W_1(Q)=-Q$
and
\[
 W_m(Q)=tW_{m-1}(BQ)+(m-1)sW_{m-2}(Q)-s^{m-1}QY^{m-1}.
\]
Then $s^mQY^m-C_m(Q)=D_YW_m(Q)$, with brackets removed
from the polynomial expression defining $C_m$.

The second identity uses the parameter connection:
\begin{equation}\label{eq:evenrec}
 uE_{j+1}(Q)=t\nabla E_j(Q)+(2j+1)E_j(Q)
                        -tE_j(gQ+Q_t).
\end{equation}
Its difference, left minus right at the polynomial level, is
$-D_Y(QY^{2j+1})$. Thus it too is an ambient identity.

Let $p_j=2j-1$ for $j\ge1$. Construct an expression $A_j(Q)\in N$
for $u^{p_j}E_j(Q)$ as follows:
\begin{align*}
 A_1(Q)&=t\nabla[Q]+[Q]-t[gQ+Q_t],\\
 A_{j+1}(Q)&=tu\nabla A_j(Q)
       +\bigl(tp_ja+(2j+1)u\bigr)A_j(Q)
       -tu A_j(gQ+Q_t).
\end{align*}
The product rule and \eqref{eq:evenrec} prove the assertion
inductively. If $V_j(Q)$ is the witness for
$u^{p_j}QY^{2j}-A_j(Q)=D_YV_j(Q)$, it is retained by
\begin{align*}
 V_1(Q)&=-QY,\\
 V_{j+1}(Q)&=tu\nabla V_j(Q)
       +\bigl(tp_ja+(2j+1)u\bigr)V_j(Q)
       -tu V_j(gQ+Q_t)-u^{p_j+1}QY^{2j+1}.
\end{align*}

Choose polynomials $\alpha_j,\beta_j\in\mathbb Z[a,t]$ with
\[
 \alpha_j u^{2j-1}+\beta_j s^{2j}=1.
\]
They exist universally, including $a=0$ and nilpotent $a$:
expand $(2u-s)^{4j-2}=1$, assigning each term with $u$ exponent
at least $2j-1$ to the first summand and every remaining term
to the second. Consequently
\begin{equation}\label{eq:evenpacket}
 E_j(Q)=\alpha_j A_j(Q)+\beta_j C_{2j}(Q)\in N.
\end{equation}
Its spatial witness is $\alpha_jV_j+\beta_jW_{2j}$.
The identity $2u-s=1$ introduces no time denominator or coefficient guard.

For example, $1=s^2+4at\,u$ supplies the first even expression.
For a time-independent base marker $Q$ it gives the concrete identity
\begin{equation}\label{eq:firsteven}
 \Phi_{QY^2}
 =4at^2\partial_t\Phi_Q+(1+2at)\Phi_Q
       +t^2\Phi_{Q(B^2-4ag)}.
\end{equation}
Its spatial witness is $-Q((1+2at)Y+tB)$: applying $D_Y$ gives
exactly the difference between the two polynomial amplitudes in
\eqref{eq:firsteven}, using $\partial_t\Phi_Q=\Phi_{PQ}$.
This polynomial identity includes $a=0$, $B=0$, and nilpotent coefficients.

\section{Odd markers and the time commutator}

Write
\[
 R_j(Q)=[QY^{2j+1}],\qquad
 L_iQ=(U_i-tg_i)Q-\partial_{U_i}Q.
\]
The two relevant spatial identities are
\begin{align}
 tR_j(BQ)&=sE_{j+1}(Q)-(2j+1)E_j(Q),\label{eq:oddstart}\\
 R_j(L_iQ)&=tE_{j+1}(B_iQ).\label{eq:oddstein}
\end{align}
Define $K_{j,F}(Q)$ to be a base-marker expression for
$tR_j(FQ)$. Equation~\eqref{eq:oddstart} starts it at $F=B$.
If it has been constructed for $F$, then
\begin{equation}\label{eq:source-derivative}
 K_{j,F_i}(Q)=K_{j,F}(L_iQ)
                          -t^2E_{j+1}(B_iFQ)
\end{equation}
constructs it for $F_i$. Indeed
$L_i(FQ)=FL_iQ-F_iQ$, and \eqref{eq:oddstein} supplies
the other term. Its spatial witness is the witness for $F$ evaluated at $L_iQ$,
minus the even witness times $t^2$, and the additional exact term
$tD_i(FQY^{2j+1})$. Hence all derivatives of $B$ are covered
without division.

Let $b_\alpha$ be the coefficient of $U^\alpha$ in $B$.
Finite Taylor translation gives the identity
\begin{equation}\label{eq:coefficient}
 b_\alpha=\frac1{\alpha!}
    \sum_{\gamma\ge0}\frac{(-U)^\gamma}{\gamma!}
                         \partial^{\alpha+\gamma}B(U).
\end{equation}
Apply each derivative expression to $(-U)^\gamma Q/\alpha!\gamma!$
and sum. This produces a base-marker expression $T_{j,\alpha}(Q)$
for $b_\alpha tR_j(Q)$ with explicit spatial witnesses.
The sum is finite and all factorials are units in every
$\mathbb Q$-algebra. Multiplication by the Taylor monomial
occurs in the amplitude before applying the expression; one must
not treat multiplication by $U_i$ as an action on the top quotient.

Finally remove the factor $t$ by the Weyl commutator, retaining
the actual parameter connection:
\begin{equation}\label{eq:oddpacket}
\begin{split}
 b_\alpha R_j(Q)={}&\nabla T_{j,\alpha}(Q)
          -T_{j,\alpha}(Q_t+gQ)-aT_{j+1,\alpha}(Q)\\
          &-b_\alpha tE_{j+1}(BQ).
\end{split}
\end{equation}
To verify it expand $\nabla(tR_j(Q))$ and use
\[
 \nabla R_j(Q)=R_j(Q_t+gQ)+aR_{j+1}(Q)+E_{j+1}(BQ).
\]
Equivalently, $[\nabla,t]=1$ removes an origin-supported
quotient after it is shown that $t$ kills every odd generator.
There is no division by $t$ and no assumption that the ambient
source is torsion-free. Applying the same combination to the
retained spatial witnesses verifies \eqref{eq:oddpacket}
before passage to any image.

\begin{samepage}
Every coefficient $b_\alpha$ therefore annihilates every odd
class modulo $N$; the even classes already belong to $N$.
This proves $I_B(M/N)=0$. If
$\sum_\alpha c_\alpha b_\alpha=1$ in $R$, summing
\eqref{eq:oddpacket} with coefficients $c_\alpha$ proves $M=N$.
This is a global coefficient-ring statement under precisely the
unimodular content hypothesis, with nilpotents and zero divisors retained.
\par\end{samepage}

The Gaussian integration-by-parts map to $R[[t]]$ is horizontal
and annihilates every displayed spatial witness. Thus all these
identities descend to the actual image. If $B=0$, every odd
period is zero coefficientwise by parity in $Y$, while the
even expression remains valid. This completes both theorems.

\section{Nilpotent coefficient content}
\label{sec:nilpotentcontent}

The coefficient condition for the actual image is weaker than the
ambient condition $I_B=R$. It also includes nonzero nilpotent sources.

\begin{corollary}\label{cor:nilpotentcontent}
If $I_B^L=0$ for some $L\ge1$, every original marked period is a
finite $D_R$-linear combination of base-marker periods. Thus $P$ is
safe exactly when all its base-marker periods are polynomial.

The same conclusion holds if a central idempotent $e\in R$ satisfies
$eI_B=eR$ and the ideal $(1-e)I_B$ is nilpotent in $(1-e)R$.
These are sufficient conditions for general $B$.
\end{corollary}
\begin{proof}
The even classes already belong to $N$ by Theorem~\ref{thm:content}.
Since $B^L=0$, the polynomial
\[
 H_L(t,U,Y)=\sum_{k=0}^{L-1}\frac{(-tBY)^k}{k!}
\]
is exactly $e^{-tBY}$. For a base marker $Q$, Gaussian parity in $Y$
gives the finite identity
\begin{equation}\label{eq:nilpotentparity}
 0=\E[QY^{2j+1}H_Le^{tP}]
   =\sum_{k=0}^{L-1}\frac{(-t)^k}{k!}
                \Phi_{QB^kY^{2j+1+k}},
\end{equation}
because $H_Le^{tP}=e^{t(aY^2+g)}$ is even in $Y$.
For every $k\ge1$, expand $B^k$ by its central coefficients, all
of which lie in $I_B$. The identity $I_B(M/N)=0$ puts each class
$[QB^kY^{2j+1+k}]$ in $N$: this uses scalar coefficient expansion,
not multiplication by $U_i$ on the quotient. Isolating the $k=0$
term in \eqref{eq:nilpotentparity} proves the assertion for every odd
marker. All sums and all operator expressions are finite.

For the last statement split $R=eR\times(1-e)R$. The unit-content
construction applies on $eR$, and the preceding nilpotent construction
applies on $(1-e)R$. Their finite expressions combine using $e$ and
$1-e$, which commute with the time operators.
\end{proof}

For example, $B^2=0$ gives the simple actual-image identity
\[
 \Phi_{QY^{2j+1}}=t\,\Phi_{QBY^{2j+2}}.
\]
Its right side is an even fibre period with a base amplitude, so the
even certificates apply directly. This includes $B=\varepsilon U_1$
over $R=\mathbb Q[\varepsilon]/(\varepsilon^2)$, where the content
ideal is nonzero and nonunit. No inverse of $a$ is required.

\section{A sufficient mixed-moment bound}

The finite expressions also give a bound on the degrees of the marked
periods.
Suppose $I_B=R$, or work over a field, and suppose the base periods
have a common bound
\[
 \deg_t\Phi_Q\le C\deg_UQ+E\qquad(Q\in R[U]),\quad C,E\ge0.
\]
For $B\ne0$ put $b=\deg B$ and
$\delta=\max(1,\deg g,2b)$. Then every original marker obeys
\begin{equation}\label{eq:cutoff}
 \deg_t\Phi_Q\le(C\delta+2)\deg Q
                    +E+C\delta(2b+3)+b+5.
\end{equation}
This bound is deliberately sufficient. It does not identify a least slope.

Here are the degree counts. For a time-independent base amplitude
of degree $q$, the even expression $E_j$ uses base amplitudes of degree
at most $q+j\delta$, parameter-derivative order at most $j$, and
polynomial time coefficients of degree at most $4j-2$ for $j\ge1$.
For $A_j$, induction gives amplitude degree $q+j\deg g$ and
time-coefficient degree $2j-1$. The conditional expression uses
degree at most $q+2jb$. The two B\'ezout factors have degrees
at most $2j-1$ and $2j-2$, giving the stated even bounds.

Each source differentiation in \eqref{eq:source-derivative} adds
at most $\max(1,\deg g-1)\le\delta$ to amplitude degree and
one to time-coefficient degree. There are at most $b$ such steps.
Taylor translation adds at most $b$ to amplitude degree.
After \eqref{eq:oddpacket}, an odd marker $Q(U)Y^{2j+1}$ therefore
uses base amplitude degrees at most $q+(j+2b+3)\delta$,
time-coefficient degrees at most $4j+b+7$, and derivative order
at most $j+2$. Since its original degree is $q+2j+1$, these
bounds imply \eqref{eq:cutoff}; differentiating a polynomial
period cannot increase its degree. A finite coefficient-content
B\'ezout identity adds no time or amplitude degree.

When $a=0$, conditional integration gives the sharper direct bound
\[
 \deg_t\Phi_{Q(U)Y^m}\le C\deg_UQ+(1+Cb)m+E.
\]
If $B=0$, all odd periods vanish and the even construction gives
a sufficient bound with $\delta=\max(1,\deg g)$; the looser
formula \eqref{eq:cutoff} may be used with $b=0$.

Each finite expression retains its base amplitudes, time operators,
spatial witnesses, and coefficient conditions. If a marker has a
nonpolynomial period, at least one of its listed base amplitudes does
too; the bounds above control the degree needed to find it.

\section{A Specialization Obstruction in Gaussian Marker Generation}
\label{sec:specialization}

Consider the one-variable Gaussian family
\[
 P_b(Y)=Y^2+bY,\qquad
 F_b(t)=\E[e^{tP_b}]
   =(1-2t)^{-1/2}\exp\!\left(\frac{b^2t^2}{2(1-2t)}\right).
\]
Here $F_b=\Phi_1$ uses the same algebraic Gaussian functional as above.
The phase is not pure, since $\E[P_b]=1$; its role is to isolate
the coefficient obstruction to marker generation.

Let $R$ again be any commutative $\mathbb Q$-algebra and let $b\in R$.
For this scalar-source family, put
\[
 V_R=\sum_{m\ge0}D_R\Phi_{Y^m}\subset R[[t]],\qquad
 N_R^{\rm act}=D_RF_b,
\]
and define the coefficient torsion ideal
\[
 \Gamma_b(R)=\{c\in R:b^Lc=0\text{ for some }L\ge1\}.
\]
The operators differentiate time; the coefficient parameter $b$
is not a differentiated variable.
Let $\delta_R$ be the point module at $t=1/2$: it is free over
$R$ on $e_j=\partial_t^j e$, $j\ge0$, with
\[
 (1-2t)e_0=0,\qquad
 (1-2t)e_j=2j e_{j-1},\qquad\partial_t e_j=e_{j+1}.
\]

\begin{theorem}\label{thm:allrings}
For every such $R,b$, the exact actual-image defect is
\begin{equation}\label{eq:defect}
 V_R/N_R^{\rm act}\simeq
     \bigl(R/(bR+\Gamma_b(R))\bigr)\otimes_R\delta_R.
\end{equation}
Consequently the unit generates the full actual module if and
only if $bR+\Gamma_b(R)=R$. This is equivalent to existence of
a finite witness
\begin{equation}\label{eq:stabilization}
 b^L(1-br)=0\quad\text{for some }L\ge0,\ r\in R.
\end{equation}
Equivalently, $R$ splits by an idempotent into a component where
$b$ is invertible and a component where $b$ is nilpotent.
\end{theorem}

Over a field this condition always holds: a nonzero $b$ is a unit,
and $b=0$ is nilpotent. Over $R=\mathbb Q[b]$, however,
$\Gamma_b(R)=0$ and the quotient is the nonzero module
$\delta_{\mathbb Q}$. The marker
\[
 \Phi_Y=\frac{bt}{1-2t}F_b
\]
does not belong to $D_RF_b$, even though the specialized marker belongs to the specialized unit-generated
module after every field specialization. Thus the nonunit-content boundary of the
base-marker theorem is a real obstruction in the actual image,
not only in an enlarged source presentation.

The same criterion fails for $\mathbb Q[[b]]$, although it holds
for every finite quotient $\mathbb Q[b]/(b^L)$. Agreement on all
finite coefficient jets does not produce a finite polynomial-Weyl
operator over the completed coefficient ring.

We first compute the quotient in the polynomial source, then identify
its exact Gaussian kernel. Removing that kernel gives the actual-image
defect. Explicit nilpotent-component operators and the specialization
sequence describe the two consequences of this calculation.

\subsection{The ambient quotient}

For this family the source and connection are
\[
 M_R=R[t,Y]/D_YR[t,Y],\qquad
 D_Y=\partial_Y-sY+tb,\qquad
 \nabla=\partial_t+Y^2+bY,\quad s=1-2t.
\]
Let $N_R=D_R[1]\subset M_R$. Theorem~\ref{thm:content}, with
no base variables, $a=1$, $B=b$ and $g=0$, gives every even
class in $N_R$ and $b(M_R/N_R)=0$. To keep the odd reduction
explicit, put $E_j=[Y^{2j}]$, $R_j=[Y^{2j+1}]$. Fibre integration
by parts gives
\begin{equation}\label{eq:wickrec}
 s[Y^m]=tb[Y^{m-1}]+(m-1)[Y^{m-2}]\qquad(m\ge1).
\end{equation}
Thus $K_j=sE_{j+1}-(2j+1)E_j=tbR_j$ lies in $N_R$, and
\begin{equation}\label{eq:content}
 bR_j=\nabla K_j-K_{j+1}-b^2tE_{j+1}\in N_R.
\end{equation}
This identity removes $t$ by the time commutator, without division.

Reduction modulo $b$ computes the entire ambient quotient.
The source commutes with this reduction because it is a cokernel
of a polynomial map. At $b=0$ the connection preserves parity.
Its even part is generated by $[1]$, and the odd part satisfies
\[
 s[Y^{2j+1}]=2j[Y^{2j-1}],\qquad
 \nabla[Y^{2j+1}]=[Y^{2j+3}].
\]
It is precisely the point module, by the inverse maps
$[Y^{2j+1}]\mapsto e_j$ and $\partial_t^je\mapsto[Y^{2j+1}]$.
Consequently
\begin{equation}\label{eq:ambientdefect}
 M_R/N_R\simeq(R/bR)\otimes_R\delta_R.
\end{equation}
In particular the ambient odd class survives when $b=0$, although
its actual Gaussian period is zero by parity. This explains why
source generation and actual-image generation have different
coefficient criteria.

\subsection{The Gaussian kernel and point torsion}

The Gaussian map $J_R:M_R\to R[[t]]$ takes $[Y^m]$ to
$\Phi_{Y^m}$. After inverting $s$, recurrence
\eqref{eq:wickrec} gives a free rank-one module with basis $[1]$.
The map there is multiplication by the invertible formal series
$F_b$. The map $R[t,s^{-1}]\to R[[t]]$ is injective: multiply by
a common power of $s$ and use injectivity of polynomials in formal
series. Writing $\Gamma_s(M_R)=\{m:s^Lm=0\text{ for some }L\ge1\}$, we obtain
\begin{equation}\label{eq:kernel}
 \ker J_R=\Gamma_s(M_R).
\end{equation}

We determine its ground classes without a flatness assumption.
The operator $D_Y$ is injective on $R[t,Y]$: for a nonzero
polynomial of highest $Y$ degree, its image has highest term
$-sY$ times that coefficient, and $s$ is a nonzerodivisor in
$R[t]$.

Suppose $s[H]=0$. There is a unique polynomial $Q$ with
$sH=D_YQ$. Reducing at $t=1/2$ gives
\[
 \partial_Y\overline Q+(b/2)\overline Q=0.
\]
Every polynomial solution is
\[
 \overline Q=c\sum_{k\ge0}\frac{(-bY/2)^k}{k!},
 \qquad c\in\Gamma_b(R),
\]
where the sum is finite on $c$. Conversely each such $c$ gives
the source ground class
\begin{equation}\label{eq:ground}
 m_c=\sum_{k=0}^{L-1}\frac{(-tb)^k}{k!}\,c[Y^{k+1}],
                  \qquad b^Lc=0.
\end{equation}
Indeed for $H_{L,c}=\sum_{k<L}(-tbY)^kc/k!$,
\begin{equation}\label{eq:torsionwitness}
 D_YH_{L,c}=-sYH_{L,c}.
\end{equation}
Changing $L$ does not change the polynomial once $b^Lc=0$.

This identifies $\ker(s:M_R\to M_R)$ with $\Gamma_b(R)$.
Here is the well-definedness and completeness check. Replacing
$H$ by $H+D_YV$ replaces $Q$ by $Q+sV$, so $\overline Q$
is intrinsic. If $\overline Q=0$, divide $Q$ by $s$; cancelling
the nonzerodivisor $s$ gives $H\in D_YR[t,Y]$. If
$\overline Q$ has constant coefficient $c$, the class $-m_c$
has the same reduction of its witness, and subtraction reduces
to the zero case. The minus sign simply reflects the witness
convention in \eqref{eq:torsionwitness}. In particular the map
$c\mapsto m_c$ is an $R$-linear isomorphism onto $\ker s$.

For clarity, we also give the elementary point-module argument
that completes this ground calculation. On any module where
$s$ acts locally nilpotently and $[\partial_t,s]=-2$, put
$D=-\partial_t/2$. Then $[D,s]=1$. The finite operator
\[
 \pi(v)=\sum_{j\ge0}\frac{D^js^jv}{j!}
\]
projects to $\ker s$. The sums are finite on each vector.
Successively applying $s$ and this projection gives the unique
expansion
\[
 v=\sum_{j\ge0}\frac{(-1)^j}{j!}D^j\pi(s^jv).
\]
The formulas follow by $sD^j=D^js-jD^{j-1}$ and telescoping;
the top surviving power proves uniqueness. Consequently
\begin{equation}\label{eq:fulltorsion}
 \Gamma_s(M_R)\simeq\Gamma_b(R)\otimes_R\delta_R,
 \quad c\otimes\partial_t^je\longmapsto\nabla^jm_c.
\end{equation}
This is the elementary point-support calculation underlying the
usual Weyl-module description; no analytic comparison or
general image-kernel computation is imported.

\subsection{The actual-image defect and the idempotent criterion}

In the ambient quotient \eqref{eq:ambientdefect}, the ground
class $m_c$ maps to $c[Y]$. Its $k\ge1$ terms in
\eqref{eq:ground} are either even, or have a coefficient
divisible by $b$, and therefore lie in $N_R$.
The image of \eqref{eq:fulltorsion} in
$(R/bR)\otimes\delta_R$ is exactly
\[
 ((\Gamma_b(R)+bR)/bR)\otimes_R\delta_R.
\]
Since $V_R=M_R/\ker J_R$ and the image of $N_R$ is
$N_R^{\rm act}$, taking the quotient proves
\eqref{eq:defect}.

The criterion \eqref{eq:stabilization} follows directly:
$1=br+c$ with $b^Lc=0$ is equivalent to
$b^L(1-br)=0$. If $L\ge1$, put $z=br$ and $e=z^L$.
Then $z^L=z^{L+1}$, so $e^2=e$, $b(re)=e$, and
$b^L(1-e)=0$. Thus $b$ is invertible on $eR$ and nilpotent
on $(1-e)R$. For $L=0$ it is already invertible. The converse
follows by combining the inverse and nilpotence witnesses on
these two components. This proves all equivalences without
Noetherian or reducedness hypotheses.

\subsection{Finite operators on a nilpotent component}

The nilpotent branch has an explicit operator, not just an
existence argument. Suppose $b^L=0$ for $L\ge1$ and put
$\nu=\lceil L/2\rceil$. Write $h=s^{-1/2}$ and define the
polynomial unit
\[
 A_L(s)=\sum_{k=0}^{\nu-1}\frac{(b^2(2-s)/8)^k}{k!}.
\]
The exact Gaussian formula gives
\[
 A_LF_b=h\sum_{j=0}^{\nu-1}\frac{b^{2j}}{8^jj!}s^{-j},
\]
because $b^2t^2/(2s)=b^2/(8s)-b^2/4+b^2s/8$.
For $E=s\partial_s=-(s/2)\partial_t$, the finite projector
\begin{equation}\label{eq:projector}
 \Pi_\nu(E)=\prod_{j=1}^{\nu-1}\frac{E+j+1/2}{j}
 \qquad\text{satisfies}\qquad
 \Pi_\nu(E)(A_LF_b)=h.
\end{equation}
Each term $s^{-j-1/2}$ is an Euler eigenvector, so the product
kills every $j\ge1$ and fixes $j=0$. Its parameter-derivative
order is at most $\nu-1$; an empty product is one.

Moreover
$\partial_t^kh=(2k-1)!!s^{-k}h$, with the $k=0$ factor one.
These rational constants are units. Conditional Wick expansion
puts every $\Phi_{Y^m}$ in $R[t,s^{-1}]h$, while
\eqref{eq:projector} puts $h$ in $D_RF_b$. Hence
\[
 V_R=D_RF_b=R[t,s^{-1}]h\quad\text{if }b^L=0.
\]
On a unit component, \eqref{eq:content} and the inverse of $b$
generate every odd class; the even classes are already generated.
For \eqref{eq:stabilization}, the idempotent $e$ combines this
construction with the finite nilpotent-component operator
\eqref{eq:projector}. Both constructions therefore give explicit
finite operators under their stated coefficient hypotheses.

\subsection{Specialization failure despite flatness}

Take $R_0=\mathbb Q[b]$. Here $\Gamma_b(R_0)=0$, so
\eqref{eq:kernel} and \eqref{eq:fulltorsion} make the Gaussian
source map injective. Therefore $V_{R_0}\simeq M_{R_0}$ and
\[
 V_{R_0}/D_{R_0}F_b\simeq\delta_{\mathbb Q}\ne0.
\]
This proves the nonmembership assertion for $\Phi_Y$; its
ambient class maps to the nonzero generator of the quotient.
There is also a concrete divisibility warning: every coefficient
of $\Phi_Y$ is divisible by $b$, yet $\Phi_Y\notin bV_{R_0}$.
Indeed $[Y]$ survives in $M_{R_0}/bM_{R_0}$ as the nonzero
odd point class. Coefficientwise divisibility in formal series
is weaker than divisibility inside the actual image.
The universal actual module is even flat over $\mathbb Q[b]$:
it is torsion-free inside the domain $\mathbb Q[b][[t]]$, and
a torsion-free module over a PID is a filtered union of finite
free submodules. This classical flatness fact is also recorded
in the Stacks Project~\cite[Lemma~15.22.11]{Stacks}.

Nevertheless the actual embedded image does not commute with
the specialization $b=0$. The source does commute with this
cokernel base change, and at $b=0$ its odd point module is killed
by the Gaussian functional. Thus the specialization map has
the exact sequence
\[
 0\longrightarrow\delta_{\mathbb Q}
 \longrightarrow V_{R_0}\otimes_{R_0}R_0/(b)
 \longrightarrow V_{b=0}\longrightarrow0.
\]
Thus generation after every field specialization and every finite
quotient $\mathbb Q[b]/(b^L)$ does not detect the universal defect.
Theorem~\ref{thm:content} cannot therefore be extended to unrestricted
coefficient rings merely by checking their field specializations.


\section{The remaining two-base question}

For $P=aY^2+B(U_1,U_2)Y+g(U_1,U_2)$, Theorem~\ref{thm:actual}
reduces purity-to-safety to one question: does $\Phi_1=1$ force
every base period

\[
 \Phi_Q=(1-2at)^{-1/2}\E_U\!\left[
 Q\exp\!\left(tg+\frac{t^2B^2}{2(1-2at)}\right)\right]
\]
to be polynomial? Wilson's theorem~\cite{Wilson} concerns powers of a single fixed
binary polynomial. The displayed parameter-dependent phase does
not satisfy that hypothesis merely because its unit series is one.
The specialization defect in Theorem~\ref{thm:allrings} concerns
marker generation for a nonpure family, so it does not decide this
remaining Gaussian moment problem.

\begin{thebibliography}{9}
\bibitem{Wilson}
M.~Wilson, \emph{A face-isolation proof of the two-variable Gaussian
Moments Conjecture}, public preprint (2026), Theorem~1,
\href{https://arxiv.org/abs/2607.23887v1}{arXiv:2607.23887v1}.

\bibitem{Stacks}
The Stacks Project Authors, \emph{The Stacks Project},
Section~15.22, Lemma~15.22.11,
\href{https://stacks.math.columbia.edu/tag/0AUW}{Tag 0AUW}.
\end{thebibliography}
\end{document}
